\section{Security Definitions and Composition (List-Based IOP Refactor)}
\label{sec:iop_refactor_security_composition}

\providecommand{\defeq}{\mathrel{:=}}
\providecommand{\pSpec}{\mathsf{pSpec}}
\providecommand{\PtoV}{\mathsf{P\_to\_V}}
\providecommand{\VtoP}{\mathsf{V\_to\_P}}
\providecommand{\Transcript}{\mathsf{Transcript}}
\providecommand{\Option}{\mathsf{Option}}
\providecommand{\StmtIn}{\mathsf{StmtIn}}
\providecommand{\StmtOut}{\mathsf{StmtOut}}
\providecommand{\StmtMid}{\mathsf{StmtMid}}
\providecommand{\WitIn}{\mathsf{WitIn}}
\providecommand{\WitOut}{\mathsf{WitOut}}
\providecommand{\WitMid}{\mathsf{WitMid}}

This section gives a self-contained, paper-style development of the security notions used in the
list-based IOP refactor (protocol specifications as lists, transcripts as heterogeneous vectors, and
\emph{outside} challenge sampling), and proves the core composition theorems used in the refactor:
completeness composition, and the implications from round-by-round (RBR) notions to their global
counterparts via a union bound.  We also record the standard $n$-fold corollaries for replicated
protocols.

\subsection{Protocol skeleton and transcripts}

\begin{definition}[Rounds and protocol specifications]
A \emph{round} is one of the following two forms:
\begin{enumerate}
  \item a prover-to-verifier message round $(\PtoV, T)$, where $T$ is the message type;
  \item a verifier-to-prover challenge round $(\VtoP, T)$, where $T$ is the challenge type.
\end{enumerate}
A \emph{protocol specification} is a finite list
\[
  \pSpec = [r_0, r_1, \dots, r_{L-1}],
\]
where each $r_k$ is a round. We write $L \defeq |\pSpec|$ for its length.
\end{definition}

\begin{definition}[Transcripts and partial transcripts]
For a protocol specification $\pSpec = [r_0,\dots,r_{L-1}]$, a \emph{full transcript} is a tuple
\[
  \tau = (\tau_0,\dots,\tau_{L-1}),
\]
where $\tau_k$ has the type specified by $r_k$ (message type if $r_k$ is a message round, challenge
type if $r_k$ is a challenge round). For $0 \le k \le L$, we write $\tau_{\le k}$ for the length-$k$
prefix $(\tau_0,\dots,\tau_{k-1})$, called a \emph{partial transcript}.
\end{definition}

\begin{definition}[Challenge positions]
Let $\mathcal{I}(\pSpec) \subseteq \{0,1,\dots,L-1\}$ denote the set of indices $i$ such that $r_i$ is
a challenge round. We call $\mathcal{I}(\pSpec)$ the set of \emph{challenge positions}.
\end{definition}

\subsection{Outside challenge sampling and the joint experiment}

\begin{definition}[Outside challenge sampling]
Fix a protocol specification $\pSpec$.  A \emph{challenge sampler} for $\pSpec$ is a randomized
procedure that samples a value $c_i$ of the appropriate challenge type for each challenge round
index $i \in \mathcal{I}(\pSpec)$, and returns the collection
\[
  \mathbf{c} = (c_i)_{i \in \mathcal{I}(\pSpec)}.
\]
We write $\mathbf{c} \gets \mathsf{Chal}(\pSpec)$ for this sampling step.
\end{definition}

\begin{remark}[Single probability space]
All security notions below are expressed as a \emph{single} probability $\Pr[\cdot]$ over the joint
randomness of:
\begin{enumerate}
  \item outside-sampled challenges $\mathbf{c} \gets \mathsf{Chal}(\pSpec)$, and
  \item any additional randomness used by the algorithms (e.g.\ shared-oracle randomness, or fresh
  randomness used when re-running the verifier on a fixed transcript).
\end{enumerate}
This is strictly weaker (and the intended notion) than a pointwise guarantee ``for all fixed
challenge strings~$\mathbf{c}$''.
\end{remark}

\subsection{Reductions, verifiers, and execution}

\begin{definition}[Plain verifier]
Fix types $\StmtIn,\StmtOut$ and a protocol specification $\pSpec$. A \emph{verifier} is a (possibly
probabilistic) algorithm
\[
  V : \StmtIn \times \Transcript(\pSpec) \to \Option(\StmtOut),
\]
where returning $\mathsf{none}$ means \emph{reject} and returning $\mathsf{some}(s)$ means
\emph{accept} with output statement $s \in \StmtOut$.
\end{definition}

\begin{definition}[Plain prover and reduction]
Fix types $\StmtIn,\WitIn,\StmtOut,\WitOut$ and $\pSpec$. A \emph{prover} is an algorithm that, given
an input pair $(x,w) \in \StmtIn \times \WitIn$ and outside-sampled challenges $\mathbf{c}$, produces
a transcript $\tau$ and an output pair $(x^P_{\mathsf{out}}, w_{\mathsf{out}}) \in \StmtOut \times
\WitOut$. A \emph{reduction} is a pair $R = (P,V)$ of a prover and verifier sharing the same
specification~$\pSpec$.
\end{definition}

\begin{definition}[Execution of a reduction]
Fix a reduction $R=(P,V)$ over $\pSpec$. On input $(x,w) \in \StmtIn\times\WitIn$, the execution
experiment samples $\mathbf{c} \gets \mathsf{Chal}(\pSpec)$, runs
\[
  (\tau, x^P_{\mathsf{out}}, w_{\mathsf{out}}) \gets P(x,w;\mathbf{c}),
\]
then computes
\[
  x^V_{\mathsf{out}} \gets V(x,\tau) \in \Option(\StmtOut).
\]
The output of the experiment is the triple
\[
  \Bigl(x^V_{\mathsf{out}},\, x^P_{\mathsf{out}},\, w_{\mathsf{out}}\Bigr).
\]
\end{definition}

\subsection{Shared oracle state and invariants}

In Lean, the prover and verifier may share access to a \emph{stateful oracle simulation}, modeled by
a state space $\Sigma$ with an initialization distribution $\mathsf{init}$ and an implementation
$\mathsf{impl}$ for oracle queries. To make sequential composition arguments unconditional (and
future-proof across models such as ROM, AGM, or hybrids), we parameterize security notions by an
oracle-state invariant
\[
  \mathsf{Inv} : \Sigma \to \{\mathsf{true},\mathsf{false}\}.
\]
We write $\Pr_{\sigma_0}[\cdot]$ for probabilities taken in the joint experiment (outside challenge
sampling plus all randomness arising from oracle answers and internal randomness) \emph{starting from
an initial oracle state $\sigma_0 \in \Sigma$}.

\paragraph{Standing hypotheses (when needed).}
We often assume:
\begin{enumerate}
  \item (\textbf{Init satisfies $\mathsf{Inv}$}) $\Pr_{\sigma_0\gets\mathsf{init}}[\mathsf{Inv}(\sigma_0)] = 1$.
  \item (\textbf{$\mathsf{impl}$ preserves $\mathsf{Inv}$}) starting from any $\sigma$ with $\mathsf{Inv}(\sigma)$,
  any oracle query step yields a post-state $\sigma'$ with $\mathsf{Inv}(\sigma')$.
\end{enumerate}

\subsection{Security notions}

\begin{definition}[Completeness]
Let $rel_{\mathsf{in}} \subseteq \StmtIn\times\WitIn$ and $rel_{\mathsf{out}} \subseteq
\StmtOut\times\WitOut$ be relations. A reduction $R=(P,V)$ over $\pSpec$ is \emph{complete with error
$\varepsilon \in \mathbb{R}_{\ge 0}$} from $rel_{\mathsf{in}}$ to $rel_{\mathsf{out}}$ if for all
$(x,w) \in rel_{\mathsf{in}}$ and all initial oracle states $\sigma_0$ with $\mathsf{Inv}(\sigma_0)=\mathsf{true}$,
\[
  \Pr_{\sigma_0}\Bigl[
    x^V_{\mathsf{out}} = \mathsf{some}(x^P_{\mathsf{out}})
    \ \wedge\ (x^P_{\mathsf{out}}, w_{\mathsf{out}}) \in rel_{\mathsf{out}}
  \Bigr] \ \ge\ 1-\varepsilon,
\]
where the probability is over the joint experiment of outside challenge sampling and shared oracle
randomness, starting from $\sigma_0$.
\end{definition}

\begin{definition}[Perfect completeness]
A reduction is \emph{perfectly complete} if it is complete with error $\varepsilon = 0$.
\end{definition}

\begin{definition}[Soundness]
Let $L_{\mathsf{in}} \subseteq \StmtIn$ and $L_{\mathsf{out}} \subseteq \StmtOut$ be languages. A
verifier $V$ over $\pSpec$ is \emph{sound with error $\varepsilon$} if for every prover $P$ (with
arbitrary auxiliary output type) and every input $x \notin L_{\mathsf{in}}$,
\[
  \Pr\bigl[\exists s \in L_{\mathsf{out}},\ V(x,\tau) = \mathsf{some}(s)\bigr] \ \le\ \varepsilon,
\]
where $\tau$ is the transcript generated by $P$ in the joint experiment.
\end{definition}

\subsection{Round-by-round soundness}

\begin{definition}[Acceptance predicate for Option-valued verifiers]
\label{def:acc_option}
For a verifier $V : \StmtIn \times \Transcript(\pSpec) \to \Option(\StmtOut)$ and a language
$L_{\mathsf{out}} \subseteq \StmtOut$, define the acceptance predicate
\[
  \mathsf{Acc}_{L_{\mathsf{out}}}(x,\tau) \defeq
  [\,\exists s \in L_{\mathsf{out}},\ V(x,\tau) = \mathsf{some}(s)\,].
\]
\end{definition}

\begin{definition}[State function]
Fix $V$, $L_{\mathsf{in}} \subseteq \StmtIn$, and $L_{\mathsf{out}} \subseteq \StmtOut$. A
\emph{(deterministic) state function} is a predicate
\[
  \mathsf{SF}(k,x,\tau_{\le k}) \in \{\mathsf{true},\mathsf{false}\}
\]
defined for all $0 \le k \le L$, all $x \in \StmtIn$, and all partial transcripts $\tau_{\le k}$,
such that:
\begin{enumerate}
  \item (\textbf{Initialization}) For all $x$, we have
  \[
    x \in L_{\mathsf{in}} \iff \mathsf{SF}(0,x,\langle\rangle)=\mathsf{true}.
  \]
  \item (\textbf{Message-round monotonicity}) For each $k < L$ such that $r_k$ is a \emph{message}
  round, for all $x$ and all $\tau_{\le k}$,
  \[
    \mathsf{SF}(k,x,\tau_{\le k})=\mathsf{false}
    \implies
    \forall m,\ \mathsf{SF}(k+1,x,\tau_{\le k}\mathbin{\|}m)=\mathsf{false},
  \]
  where $m$ ranges over messages of the appropriate type and $\tau_{\le k}\mathbin{\|}m$ denotes
  transcript concatenation.
  \item (\textbf{Terminal soundness}) For all $x$ and full transcripts $\tau$, if
  $\mathsf{SF}(L,x,\tau)=\mathsf{false}$ then
  \[
    \forall \sigma_0,\ \mathsf{Inv}(\sigma_0)=\mathsf{true}\ \implies\
      \Pr_{\sigma_0}\bigl[\mathsf{Acc}_{L_{\mathsf{out}}}(x,\tau)\bigr]=0.
  \]
\end{enumerate}
\end{definition}

\begin{definition}[RBR soundness]
Let $\delta : \mathcal{I}(\pSpec) \to \mathbb{R}_{\ge 0}$ be an error function on challenge
positions. We say $V$ is \emph{round-by-round sound} with error function $\delta$ if there exists a
state function $\mathsf{SF}$ such that for every $x \notin L_{\mathsf{in}}$, every prover $P$, and
every challenge position $i \in \mathcal{I}(\pSpec)$, and every initial oracle state $\sigma_0$ with
$\mathsf{Inv}(\sigma_0)=\mathsf{true}$,
\[
  \Pr_{\sigma_0}\Bigl[
    \mathsf{SF}(i,x,\tau_{\le i})=\mathsf{false}
    \ \wedge\
    \mathsf{SF}(i+1,x,\tau_{\le i+1})=\mathsf{true}
  \Bigr]
  \ \le\ \delta(i),
\]
where $\tau$ is the transcript produced by $P$ in the joint experiment starting from $\sigma_0$.
\end{definition}

\begin{lemma}[Terminal soundness eliminates acceptance on bad terminal states]
\label{lem:sf_accept_and_not_terminal_zero}
Fix $x$ and a prover $P$, and let $\tau$ be the random transcript produced in the joint experiment.
Let $A \defeq \mathsf{Acc}_{L_{\mathsf{out}}}(x,\tau)$ be the acceptance event. If $\mathsf{SF}$
satisfies terminal soundness, then
\[
  \Pr\bigl[A \ \wedge\ \mathsf{SF}(L,x,\tau)=\mathsf{false}\bigr] = 0,
\]
and therefore
\[
  \Pr[A]
  = \Pr\bigl[A \ \wedge\ \mathsf{SF}(L,x,\tau)=\mathsf{true}\bigr]
  \ \le\ \Pr\bigl[\mathsf{SF}(L,x,\tau)=\mathsf{true}\bigr].
\]
\end{lemma}

\begin{proof}
Unfold the joint experiment as ``sample $\tau$ using $P$, then run $V$ on $(x,\tau)$ using its
internal randomness''. For any fixed transcript value $\tau_0$ with $\mathsf{SF}(L,x,\tau_0)=\mathsf{false}$,
terminal soundness gives $\Pr[\mathsf{Acc}_{L_{\mathsf{out}}}(x,\tau_0)] = 0$. Averaging over the random
choice of $\tau$ implies $\Pr[A \wedge \mathsf{SF}(L,x,\tau)=\mathsf{false}] = 0$. The second display
follows by the partition $\Pr[A] = \Pr[A \wedge \mathsf{SF}(L)=\mathsf{true}] + \Pr[A \wedge
\mathsf{SF}(L)=\mathsf{false}]$.
\end{proof}

\begin{lemma}[Existence of a flip at a challenge position]
\label{lem:sf_exists_challenge_flip}
Fix $x$ and a full transcript $\tau$. Assume:
\begin{enumerate}
  \item $\mathsf{SF}(0,x,\langle\rangle)=\mathsf{false}$, and
  \item $\mathsf{SF}(L,x,\tau)=\mathsf{true}$.
\end{enumerate}
Then there exists a challenge position $i \in \mathcal{I}(\pSpec)$ such that
\[
  \mathsf{SF}(i,x,\tau_{\le i})=\mathsf{false}
  \quad\text{and}\quad
  \mathsf{SF}(i+1,x,\tau_{\le i+1})=\mathsf{true}.
\]
\end{lemma}

\begin{proof}
Let $k$ be the smallest index in $\{0,1,\dots,L\}$ such that $\mathsf{SF}(k,x,\tau_{\le k})=\mathsf{true}$.
Such a $k$ exists by assumption (2), and $k \neq 0$ by assumption (1). Hence $k=i+1$ for some
$i \in \{0,\dots,L-1\}$, and by minimality we have $\mathsf{SF}(i,x,\tau_{\le i})=\mathsf{false}$ and
$\mathsf{SF}(i+1,x,\tau_{\le i+1})=\mathsf{true}$.

It remains to show $i$ is a challenge position. If $r_i$ were a message round, message-round
monotonicity would imply that from $\mathsf{SF}(i,x,\tau_{\le i})=\mathsf{false}$ it follows that
$\mathsf{SF}(i+1,x,\tau_{\le i+1})=\mathsf{false}$, contradicting
$\mathsf{SF}(i+1,x,\tau_{\le i+1})=\mathsf{true}$. Therefore $r_i$ is a challenge round, i.e.
$i\in\mathcal{I}(\pSpec)$.
\end{proof}

\begin{theorem}[RBR soundness implies soundness via a union bound]
\label{thm:rbr_soundness_implies_soundness_list}
If $V$ is RBR-sound for $(L_{\mathsf{in}},L_{\mathsf{out}})$ with error function
$\delta : \mathcal{I}(\pSpec) \to \mathbb{R}_{\ge 0}$, then $V$ is sound for
$(L_{\mathsf{in}},L_{\mathsf{out}})$ with error
\[
  \varepsilon \defeq \sum_{i \in \mathcal{I}(\pSpec)} \delta(i).
\]
\end{theorem}

\begin{proof}
Fix any $x \notin L_{\mathsf{in}}$ and any prover $P$, and let $\tau$ be the resulting transcript in
the joint experiment. Consider the acceptance event
\[
  A \defeq \mathsf{Acc}_{L_{\mathsf{out}}}(x,\tau).
\]
By initialization, $\mathsf{SF}(0,x,\langle\rangle)=\mathsf{false}$. By terminal soundness,
\Cref{lem:sf_accept_and_not_terminal_zero} yields
\[
  \Pr[A] \le \Pr\bigl[\mathsf{SF}(L,x,\tau)=\mathsf{true}\bigr].
\]
Moreover, by message-round monotonicity, the predicate $\mathsf{SF}$ cannot change from
$\mathsf{false}$ to $\mathsf{true}$ across message rounds; hence if $\mathsf{SF}(0,\dots)$ is false
and $\mathsf{SF}(L,\dots)$ is true, there must exist a challenge position $i \in \mathcal{I}(\pSpec)$
at which a flip occurs:
\[
  \bigl[\mathsf{SF}(L,x,\tau)=\mathsf{true}\bigr]
  \subseteq
  \bigcup_{i \in \mathcal{I}(\pSpec)}
  \Bigl[\mathsf{SF}(i,x,\tau_{\le i})=\mathsf{false}\ \wedge\ \mathsf{SF}(i+1,x,\tau_{\le i+1})=\mathsf{true}\Bigr].
\]
This inclusion follows by applying \Cref{lem:sf_exists_challenge_flip} pointwise to each outcome of
$\tau$.
Taking probabilities and applying the union bound yields
\[
  \Pr[A]
  \le \Pr\bigl[\mathsf{SF}(L,x,\tau)=\mathsf{true}\bigr]
  \le \sum_{i \in \mathcal{I}(\pSpec)}
  \Pr\Bigl[\mathsf{SF}(i,x,\tau_{\le i})=\mathsf{false}\ \wedge\ \mathsf{SF}(i+1,x,\tau_{\le i+1})=\mathsf{true}\Bigr]
  \le
  \sum_{i \in \mathcal{I}(\pSpec)} \delta(i),
\]
as required.
\end{proof}

\subsection{Round-by-round knowledge soundness}

\begin{definition}[Straightline knowledge soundness]
Let $rel_{\mathsf{in}} \subseteq \StmtIn\times\WitIn$ and $rel_{\mathsf{out}} \subseteq
\StmtOut\times\WitOut$. A verifier $V$ is \emph{(straightline) knowledge sound} with error
$\varepsilon$ if there exists an extractor
\[
  E : \StmtIn \times \WitOut \times \Transcript(\pSpec) \to \WitIn
\]
such that for every prover $P$ and every input statement $x$, in the joint experiment we have
\[
  \Pr\Bigl[
    \bigl(\exists s,\ V(x,\tau)=\mathsf{some}(s)\ \wedge\ (s,w_{\mathsf{out}})\in rel_{\mathsf{out}}\bigr)
    \ \wedge\
    (x,E(x,w_{\mathsf{out}},\tau)) \notin rel_{\mathsf{in}}
  \Bigr]
  \ \le\ \varepsilon.
\]
\end{definition}

\begin{definition}[RBR extractor and intermediate witnesses]
Fix $\pSpec$ with $L=|\pSpec|$. Let $(W_k)_{k=0}^{L}$ be a family of \emph{intermediate witness}
types such that $W_0 = \WitIn$. An \emph{RBR extractor} is a pair of deterministic maps:
\begin{enumerate}
  \item (\textbf{Backward step}) For each $k \in \{0,\dots,L-1\}$, a function
  \[
    \mathsf{extMid}_k : \StmtIn \times \Transcript(\pSpec)_{\le k+1} \times W_{k+1} \to W_k.
  \]
  \item (\textbf{Finalization}) A function
  \[
    \mathsf{extOut} : \StmtIn \times \Transcript(\pSpec) \times \WitOut \to W_L.
  \]
\end{enumerate}
\end{definition}

\begin{definition}[Knowledge state function]
Fix relations $rel_{\mathsf{in}} \subseteq \StmtIn\times\WitIn$ and $rel_{\mathsf{out}} \subseteq
\StmtOut\times\WitOut$, a verifier $V$, intermediate witness types $(W_k)_{k=0}^L$ with $W_0=\WitIn$,
and an RBR extractor $(\mathsf{extMid}_k,\mathsf{extOut})$. A \emph{knowledge state function} is a
predicate
\[
  \mathsf{KSF}(k,x,\tau_{\le k},w_k) \in \{\mathsf{true},\mathsf{false}\}
\]
defined for $0 \le k \le L$, $x\in\StmtIn$, partial transcripts $\tau_{\le k}$, and
$w_k \in W_k$, such that:
\begin{enumerate}
  \item (\textbf{Initialization}) For all $x$ and $w_0 \in W_0=\WitIn$,
  \[
    (x,w_0) \in rel_{\mathsf{in}} \iff \mathsf{KSF}(0,x,\langle\rangle,w_0)=\mathsf{true}.
  \]
  \item (\textbf{Backward propagation on message rounds}) If $r_k$ is a message round, then for all
  $x$, all $\tau_{\le k}$, all messages $m$, and all $w_{k+1}\in W_{k+1}$,
  \[
    \mathsf{KSF}(k+1,x,\tau_{\le k}\mathbin{\|}m,w_{k+1})=\mathsf{true}
    \implies
    \mathsf{KSF}(k,x,\tau_{\le k},\mathsf{extMid}_k(x,\tau_{\le k}\mathbin{\|}m,w_{k+1}))=\mathsf{true}.
  \]
  \item (\textbf{Finalization}) For all $x$, full transcripts $\tau$, output witnesses
  $w_{\mathsf{out}}\in\WitOut$, and initial oracle states $\sigma_0$ with
  $\mathsf{Inv}(\sigma_0)=\mathsf{true}$, if
  \[
    \Pr_{\sigma_0}\Bigl[\exists s,\ V(x,\tau) = \mathsf{some}(s)\ \wedge\ (s,w_{\mathsf{out}})\in rel_{\mathsf{out}}\Bigr] > 0,
  \]
  then
  \[
    \mathsf{KSF}(L,x,\tau,\mathsf{extOut}(x,\tau,w_{\mathsf{out}}))=\mathsf{true}.
  \]
\end{enumerate}
\end{definition}

\begin{definition}[RBR knowledge soundness]
Let $\kappa : \mathcal{I}(\pSpec)\to \mathbb{R}_{\ge 0}$ be an error function. We say that $(V,
\mathsf{KSF})$ satisfies \emph{RBR knowledge soundness} with error function $\kappa$ if for every
statement $x$, every prover $P$, every challenge position $i \in \mathcal{I}(\pSpec)$, and every
initial oracle state $\sigma_0$ with $\mathsf{Inv}(\sigma_0)=\mathsf{true}$,
\[
  \Pr_{\sigma_0}\Bigl[
    \exists w_{i+1}\in W_{i+1},\
    \mathsf{KSF}(i+1,x,\tau_{\le i+1},w_{i+1})=\mathsf{true}
    \ \wedge\
    \mathsf{KSF}(i,x,\tau_{\le i},\mathsf{extMid}_i(x,\tau_{\le i+1},w_{i+1}))=\mathsf{false}
  \Bigr]
  \ \le\ \kappa(i),
\]
where $\tau$ is the transcript produced by $P$ in the joint experiment starting from $\sigma_0$.
\end{definition}

\begin{theorem}[RBR knowledge soundness implies knowledge soundness]
\label{thm:rbr_ks_implies_ks_list}
Assume there exist intermediate witness types $(W_k)_{k=0}^L$ with $W_0=\WitIn$, an RBR extractor
$(\mathsf{extMid}_k,\mathsf{extOut})$, and a knowledge state function $\mathsf{KSF}$ such that $(V,
\mathsf{KSF})$ satisfies RBR knowledge soundness with error function
$\kappa:\mathcal{I}(\pSpec)\to\mathbb{R}_{\ge 0}$. Define the straightline extractor $E$ by:
\begin{align*}
  w_L &\defeq \mathsf{extOut}(x,\tau,w_{\mathsf{out}}),\\
  w_k &\defeq \mathsf{extMid}_k(x,\tau_{\le k+1},w_{k+1}) \qquad (k=L-1,\dots,0),\\
  E(x,w_{\mathsf{out}},\tau) &\defeq w_0.
\end{align*}
Then $V$ is knowledge sound for $(rel_{\mathsf{in}},rel_{\mathsf{out}})$ with error
\[
  \varepsilon \defeq \sum_{i\in \mathcal{I}(\pSpec)} \kappa(i).
\]
\end{theorem}

\begin{proof}
Fix any statement $x$ and any prover $P$, and let $(\tau, w_{\mathsf{out}})$ be the transcript and
output witness produced by $P$ in the joint experiment. Consider the ``bad'' event
\[
  B \defeq
  \Bigl[
    \bigl(\exists s,\ V(x,\tau)=\mathsf{some}(s)\ \wedge\ (s,w_{\mathsf{out}})\in rel_{\mathsf{out}}\bigr)
    \ \wedge\
    (x,E(x,w_{\mathsf{out}},\tau)) \notin rel_{\mathsf{in}}
  \Bigr].
\]
Unfolding $E$, the second conjunct is $(x,w_0)\notin rel_{\mathsf{in}}$. By initialization of
$\mathsf{KSF}$, this implies $\mathsf{KSF}(0,x,\langle\rangle,w_0)=\mathsf{false}$.

Define the \emph{acceptance probability} for the fixed pair $(\tau,w_{\mathsf{out}})$ by
\[
  p_{\mathsf{out}}(\tau,w_{\mathsf{out}})
  \defeq
  \Pr\Bigl[\exists s,\ V(x,\tau) = \mathsf{some}(s)\ \wedge\ (s,w_{\mathsf{out}})\in rel_{\mathsf{out}}\Bigr],
\]
where the probability is over the (fresh) randomness used by $V$ when evaluated on the fixed
transcript $\tau$ (and all other randomness is held fixed).

If $p_{\mathsf{out}}(\tau,w_{\mathsf{out}})=0$, then the event inside the probability has probability
$0$ and hence cannot contribute to $\Pr[B]$. Therefore,
\[
  \Pr[B] = \Pr\bigl[B \ \wedge\ p_{\mathsf{out}}(\tau,w_{\mathsf{out}}) > 0\bigr].
\]
On the event $p_{\mathsf{out}}(\tau,w_{\mathsf{out}}) > 0$, the finalization axiom implies
\[
  \mathsf{KSF}(L,x,\tau,w_L)=\mathsf{true}
  \qquad\text{where } w_L=\mathsf{extOut}(x,\tau,w_{\mathsf{out}}).
\]

Now consider the backward-propagated witnesses $(w_k)_{k=0}^L$ induced by the RBR extractor along the
prefixes $\tau_{\le k}$. By backward propagation on message rounds, whenever $r_k$ is a message
round, truth of $\mathsf{KSF}(k+1,x,\tau_{\le k+1},w_{k+1})$ implies truth of
$\mathsf{KSF}(k,x,\tau_{\le k},w_k)$. Therefore, the predicate $\mathsf{KSF}$ can only change from
false to true (when moving forward in $k$) across \emph{challenge} positions.

Since $\mathsf{KSF}(0,\dots)=\mathsf{false}$ but $\mathsf{KSF}(L,\dots)=\mathsf{true}$ on $B$, there
must exist $i\in\mathcal{I}(\pSpec)$ such that
\[
  \mathsf{KSF}(i+1,x,\tau_{\le i+1},w_{i+1})=\mathsf{true}
  \quad\text{and}\quad
  \mathsf{KSF}(i,x,\tau_{\le i},w_i)=\mathsf{false}.
\]
By definition of $w_i$, this is precisely the RBR-KS ``flip'' event at~$i$ witnessed by
$w_{i+1}$. Hence we have the inclusion
\[
  \bigl[B \ \wedge\ p_{\mathsf{out}}(\tau,w_{\mathsf{out}}) > 0\bigr]
  \subseteq \bigcup_{i\in\mathcal{I}(\pSpec)} F_i,
\]
where each $F_i$ is the existential flip event in the RBR knowledge soundness definition. Applying
the union bound and the assumed per-round bounds gives
\[
  \Pr[B]
  = \Pr\bigl[B \ \wedge\ p_{\mathsf{out}}(\tau,w_{\mathsf{out}}) > 0\bigr]
  \le \sum_{i\in\mathcal{I}(\pSpec)} \Pr[F_i]
  \le \sum_{i\in\mathcal{I}(\pSpec)} \kappa(i),
\]
as desired.
\end{proof}

\subsection{Completeness under sequential composition}

\begin{definition}[Sequential composition of reductions]
Let $R_1=(P_1,V_1)$ be a reduction over $\pSpec_1$ from $(\StmtIn_1,\WitIn_1)$ to $(\StmtMid,\WitMid)$,
and let $R_2=(P_2,V_2)$ be a reduction over $\pSpec_2$ from $(\StmtMid,\WitMid)$ to $(\StmtOut,\WitOut)$.
Their \emph{sequential composition} $R = R_2 \circ R_1$ is the reduction over $\pSpec_1 \mathbin{++}
\pSpec_2$ defined by: on input $(x,w)$, run $P_1$ to obtain an intermediate output $(y,u)$ and prefix
transcript $\tau^{(1)}$, then run $P_2$ on $(y,u)$ to obtain $(z,v)$ and transcript $\tau^{(2)}$;
output transcript $\tau \defeq \tau^{(1)} \mathbin{\|} \tau^{(2)}$, prover output $(z,v)$, and accept
only if both verifiers accept in sequence (equivalently, run $V_1$ on $\tau^{(1)}$ to obtain $y$ and
then run $V_2$ on $\tau^{(2)}$ with input $y$).
\end{definition}

\begin{theorem}[Completeness composes with additive error]
\label{thm:completeness_comp_add}
Let $rel_{\mathsf{in}} \subseteq \StmtIn_1\times\WitIn_1$, $rel_{\mathsf{mid}} \subseteq
\StmtMid\times\WitMid$, and $rel_{\mathsf{out}} \subseteq \StmtOut\times\WitOut$. If $R_1$ is
complete from $rel_{\mathsf{in}}$ to $rel_{\mathsf{mid}}$ with error $\varepsilon_1$ and $R_2$ is
complete from $rel_{\mathsf{mid}}$ to $rel_{\mathsf{out}}$ with error $\varepsilon_2$, then the
composition $R_2\circ R_1$ is complete from $rel_{\mathsf{in}}$ to $rel_{\mathsf{out}}$ with error
$\varepsilon_1+\varepsilon_2$.
\end{theorem}

\begin{proof}
Fix any $(x,w)\in rel_{\mathsf{in}}$. Let $G_1$ be the event that $R_1$'s verifier accepts with
output equal to $R_1$'s prover statement and $(y,u)\in rel_{\mathsf{mid}}$; by completeness of $R_1$,
$\Pr[\neg G_1]\le \varepsilon_1$.  Conditional on $G_1$ (and thus on $(y,u)\in rel_{\mathsf{mid}}$),
let $G_2$ be the analogous ``good'' event for $R_2$ with respect to $rel_{\mathsf{out}}$; by
completeness of $R_2$ applied to the specific valid intermediate pair, we have
\[
  \Pr[\neg G_2 \mid G_1] \le \varepsilon_2.
\]
The composed reduction succeeds exactly on $G_1 \wedge G_2$, hence
\[
  \Pr[G_1 \wedge G_2]
  \ge 1 - \Pr[\neg G_1] - \Pr[\neg G_2 \wedge G_1]
  \ge 1 - \varepsilon_1 - \varepsilon_2,
\]
where the last inequality uses $\Pr[\neg G_2 \wedge G_1] = \Pr[G_1]\Pr[\neg G_2\mid G_1] \le
\varepsilon_2$. This is completeness with error $\varepsilon_1+\varepsilon_2$.
\end{proof}

\begin{corollary}[$n$-fold completeness]
\label{cor:completeness_comp_n}
If a reduction is complete from $rel$ to $rel$ with error $\varepsilon$, then the $n$-fold sequential
composition of the reduction with itself is complete from $rel$ to $rel$ with error $n\varepsilon$.
\end{corollary}

\begin{proof}
By induction on $n$, applying \Cref{thm:completeness_comp_add} at each step.
\end{proof}

\begin{corollary}[$n$-fold soundness bounds from RBR bounds]
\label{cor:rbr_soundness_comp_n}
Let $\pSpec^{\mathsf{rep}(n)}$ denote the replicated protocol that concatenates $n$ copies of
$\pSpec$. If each copy admits an RBR bound $\delta:\mathcal{I}(\pSpec)\to\mathbb{R}_{\ge0}$, then the
replicated verifier is sound with error at most
\[
  n \cdot \sum_{i\in\mathcal{I}(\pSpec)} \delta(i).
\]
An analogous statement holds for knowledge soundness with $\kappa$ in place of $\delta$.
\end{corollary}

\begin{proof}
The replicated protocol has $n$ disjoint blocks of challenge positions, each block contributing a
union-bound term $\sum_i \delta(i)$ (resp.\ $\sum_i \kappa(i)$). Summing over blocks gives the stated
bound.
\end{proof}

